\chapter{Angles et cercle trigonométrique}
\label{workbook:chapter:14}
\section{Degrés et radians}
\label{workbook:section:14.01}
\begin{exerciseblock}{Degrés et radians}
\item Convertir en radians : $30^\circ$, $135^\circ$, $-210^\circ$, $450^\circ$.
\item Convertir en degrés : $\pi/8$, $5\pi/6$, $-7\pi/4$, $3$ radians.
\item Expliquer pourquoi la conversion utilise la proportion $180^\circ=\pi$ radians.
\end{exerciseblock}
\section{Le cercle unité}
\label{workbook:section:14.02}
\begin{exerciseblock}{Le cercle unité}
\item Placer sur le cercle unité les angles $0$, $\pi/2$, $\pi$, $3\pi/2$ et donner $(\cos\theta,\sin\theta)$.
\item Pour un point $P(-3/5,4/5)$ sur le cercle unité, déterminer sinus, cosinus et quadrant.
\item Vérifier l'identité $\cos^2\theta+\sin^2\theta=1$ à partir de l'équation du cercle.
\end{exerciseblock}
\section{Angles remarquables}
\label{workbook:section:14.03}
\begin{exerciseblock}{Angles remarquables}
\item Construire les valeurs exactes de sinus et cosinus pour $30^\circ$, $45^\circ$ et $60^\circ$ à partir des triangles remarquables.
\item Calculer exactement $\sin(150^\circ)$, $\cos(225^\circ)$ et $\tan(300^\circ)$.
\item Compléter un tableau des signes de sinus, cosinus et tangente dans les quatre quadrants.
\end{exerciseblock}
\section{Périodicité et symétries}
\label{workbook:section:14.04}
\begin{exerciseblock}{Périodicité et symétries}
\item Réduire les angles $17\pi/6$, $-13\pi/4$ et $25\pi/3$ modulo $2\pi$.
\item Utiliser les symétries pour calculer $\sin(-\pi/3)$, $\cos(-\pi/3)$ et $\tan(\pi+\pi/6)$.
\item Justifier que sinus est impaire, cosinus paire et tangente de période $\pi$.
\end{exerciseblock}
\section{Le radian mesure naturellement une rotation}
\label{workbook:section:14.05}
\begin{exerciseblock}{Le radian mesure naturellement une rotation}
\item Sur un cercle de rayon $5$, calculer la longueur d'arc associée à $1{,}2$ rad.
\item Un arc mesure $18$ cm sur un cercle de rayon $12$ cm. Trouver l'angle en radians puis en degrés.
\item Comparer le sens géométrique de $1$ radian sur des cercles de rayons différents.
\end{exerciseblock}
\section{Pourquoi les radians simplifient les formules}
\label{workbook:section:14.06}
\begin{exerciseblock}{Pourquoi les radians simplifient les formules}
\item Comparer $s=r\theta$ avec une formule utilisant les degrés et expliquer la simplicité du radian.
\item Pour une petite rotation $\theta=0{,}02$ rad sur un cercle de rayon $3$, estimer le déplacement le long de l'arc.
\item Expliquer pourquoi la vitesse angulaire en rad/s se combine directement avec le temps dans $\theta=\omega t$.
\end{exerciseblock}
\section{Mesurer et lire une rotation}
\label{workbook:section:14.07}
\begin{exerciseblock}{Mesurer et lire une rotation}
\item Une roue tourne de $7\pi/3$ radians. Donner le nombre de tours et l'angle final principal.
\item Calculer la longueur d'arc et l'aire du secteur pour $r=4$ et $\theta=5\pi/6$.
\item Une calculatrice en mode degrés donne $\sin(\pi/6)\approx0{,}0091$. Diagnostiquer l'erreur.
\end{exerciseblock}
\section*{Synthèse du chapitre}
\begin{synthesisbox}
\item Un engrenage de rayon $8$ cm tourne de $135^\circ$. Calculer l'arc parcouru par un point du bord et l'aire balayée.
\item Déterminer toutes les coordonnées exactes des points du cercle unité dont l'ordonnée vaut $\sqrt3/2$.
\item Une rotation de $-11\pi/6$ et une rotation de $\pi/6$ conduisent-elles au même point? Expliquer avec angles coterminaux et orientation.
\end{synthesisbox}
